\documentclass[10pt,a4paper]{amsart}
\usepackage[margin=19mm]{geometry}
\usepackage{amsmath,amssymb,mathtools}
\usepackage[scr=rsfs,cal=euler]{mathalfa}
\usepackage{xfrac}
\usepackage[unicode,hidelinks,bookmarks=false]{hyperref}
\newcommand{\ie}{i.e.\ }
\newcommand{\cf}{cf.\ }
\newcommand{\resp}{respectively\ }
\newcommand{\Subsection}{\subsection}
\providecommand{\nobreakdash}{\nobreak\hbox{-}}
\setlength{\emergencystretch}{2em}
\multlinegap0pt
\usepackage{amssymb}
\usepackage{bm}
\usepackage{mathtools}
\newtheorem{theorem}{Theorem}
\newtheorem{proposition}{Proposition}
\title{\bf A Guide to Applications of $k$-Contact Geometry\\ in Dissipative Field Equations}
\author{J. de Lucas, J. Lange,, M. Krych}
\date{Source: arXiv:2605.13313v1 (CC BY 4.0)}
\begin{document}
\maketitle

\noindent\emph{Source and licence.} \bf A Guide to Applications of $k$-Contact Geometry\\ in Dissipative Field Equations. Version arXiv:2605.13313v1 (CC BY 4.0), distributed by arXiv under the Creative Commons Attribution 4.0 International licence, http://creativecommons.org/licenses/by/4.0/, as recorded in the arXiv licence metadata for this exact version and shown on the abstract page https://arxiv.org/abs/2605.13313v1. Full licence notice: \texttt{licenses/arxiv-2605.13313-CC-BY-4.0.txt}.\par\smallskip

\noindent\textit{Editorial scope.} Counted unit: the article's proposition prop:regularity\_not\_hyperbolicity (source0.tex lines 573--577). The article's complete original proof of it is delivered with it (source0.tex lines 579--607, one \texttt{proof} environment of 29 lines). No mathematics is written, repaired or re-derived: the statement and the proof are the article's own lines, in the article's own notation, and no retained line is substituted or re-flowed.  Retained uncounted as the source-local context the proof needs: lines 556--564.  Nothing was omitted from any retained span, so the package carries no omission record.  No retained line contains a citation, so the wrapper carries no bibliography and no bibliography entry.  No retained line contains a figure environment, an image-inclusion command or drawing code, so no image is delivered and the package declares no asset dependency.  Editorial formatting only, in addition: an AMS \texttt{amsart} preamble that restates the article's own declarations and macros used by the retained lines, namely the environments \texttt{theorem}, \texttt{proposition} on separate counters, together with the standard packages those lines need.  The retained environments print in this document as Theorem 1, Proposition 1 in order of appearance, whereas the article prints the same items with the numbering of its own full document.  The complete native source of the article as distributed in the arXiv e-print for this exact version is bundled unchanged as \texttt{sources/200494/source0.tex}.
\begin{theorem}\label{thm:linear_z_same_second_order_system}
Let \(h\in C^\infty(\mathcal J_{Q,k})\) take the form \(h(q^i,p_i^\alpha,z^\alpha)=g(q^i,p_i^\alpha)+\sum_{\alpha=1}^kA_\alpha (z^\alpha)^{\lambda_\alpha}\), where \(A_1,\dots,A_k\) are real constants, $\lambda_1,\ldots,\lambda_k$ are non-negative integers, and \(h\) is regular. Let  \(\psi(t)=(q^i(t),p_i^\alpha(t),z^\alpha(t))\) be an integral section of an \(\bm\eta_{Q,k}\)-Hamiltonian \(k\)-vector field \({\bm X}_{h}\). Then, \(\psi(t)\) projects onto a map \(q(t)=(q^1(t),\dots,q^n(t))\) satisfying a second-order system of PDEs. More precisely, if \(p_i^\alpha=P_i^\alpha(q ,v_\beta^j)\) denotes the local inverse of the fibre derivative, determined by \(v_\beta^i=\partial h/\partial p_i^\beta=\partial g/\partial p_i^\beta\), then the functions \(q^i(t)\) satisfy the second-order system of PDEs of the form
\begin{equation}\label{eq:second_order_system_linear_z}
\sum_{\alpha=1}^k \frac{\partial}{\partial t^\alpha}\Bigl(P_i^\alpha(q,v_\beta)\Bigr)
+\frac{\partial g}{\partial q^i}\Bigl(q,P_j^\gamma(q,v_\beta)\Bigr)
+\sum_{\alpha=1}^k R(z_\alpha,\lambda_\alpha) A_\alpha P_i^\alpha(q,v_\beta)=0,\qquad i=1,\ldots,n\,,
\end{equation}
where \(q=q(t)\) and \(v_\beta =\partial q /\partial t^\beta\) and $R(z ,\mu)=z ^{\mu-1}\mu$  for $\mu\in \mathbb{N}_+$ and $R(z ,0)=0$ for every $z\in \mathbb{R}$.
\end{theorem}

\begin{proposition}\label{prop:regularity_not_hyperbolicity}
Let \((J_{Q,k},\bm\eta_{Q,k},h)\) be a canonical \(k\)-contact Hamiltonian system with $\dim Q=1$, and assume  the associated  second-order system of PDEs as in Theorem~\ref{thm:linear_z_same_second_order_system}.   
If $A=\left(\frac{\partial^2 h}{\partial p^\alpha\partial p^\beta}\right)$ is definite, the resulting second-order equation is elliptic. If $A$ has signature $(k-1,1)$ or $(1,k-1)$, one obtains a hyperbolic equation. If $A$ is indefinite with at least two positive and two negative eigenvalues, the equation is of ultrahyperbolic type.  

\end{proposition} 

\begin{proof}
Since \(\dim Q=1\), the reconstructed second-order system reduces to a single second-order equation. Hence its analytic type is determined by the quadratic form associated with its principal symbol. By Theorem~\ref{thm:linear_z_same_second_order_system}, if \(v_\beta=\partial q/\partial t^\beta\), the reconstructed equation has the form
\[
\sum_{\alpha=1}^k \frac{\partial}{\partial t^\alpha}\Bigl(P^\alpha(q,v_\beta)\Bigr)
+\frac{\partial g}{\partial q}\Bigl(q,P^\gamma(q,v_\beta)\Bigr)
+\sum_{\alpha=1}^k R(z^\alpha,\lambda_\alpha)A_\alpha P^\alpha(q,v_\beta)=0.
\]
Therefore, the principal part of the equation, namely the one concerning second-order derivatives of the unknown variables, is
\[
\sum_{\alpha,\beta=1}^k \frac{\partial P^\alpha}{\partial v_\beta}(q,v_\gamma)\,
\frac{\partial^2 q}{\partial t^\alpha\partial t^\beta},
\]
and its principal symbol is the quadratic form
\[
\sigma(\xi)=\sum_{\alpha,\beta=1}^k \frac{\partial P^\alpha}{\partial v_\beta}(q,v_\gamma)\,\xi_\alpha\xi_\beta.
\]

Now, the regularity of \(h\) implies that the Hessian matrix
$A$
is invertible. Since \(P^\alpha\) is obtained locally by inverting the relations \(v_\alpha=\partial h/\partial p^\alpha\), the matrix
\(
\left(\frac{\partial P^\alpha}{\partial v_\beta}\right)
\)
is precisely \(A^{-1}\). Hence the principal symbol is determined by the quadratic form associated with \(A^{-1}\), and therefore by the index of \(A\), since \(A\) and \(A^{-1}\) have the same inertia.

Consequently, if \(A\) is definite, then \(\sigma\) is definite and the equation is elliptic. If \(A\) has signature \((k-1,1)\) or \((1,k-1)\), then \(\sigma\) has the same signature and the equation is hyperbolic. More generally, if \(A\) is indefinite with at least two positive and two negative eigenvalues, then the equation is of ultrahyperbolic type.  


\end{proof}

\end{document}
